% !TEX root = main.tex
% 由 前期项目研究总结.md 转换；保持严谨求实的工科语言风格与真实实验数据。

\begin{center}
{\zihao{3}\bfseries GearPro 齿轮外观缺陷检测系统：前期项目研究总结\par}
\end{center}

\section*{一、 项目选题的背景}

齿轮是机械动力传动中核心的基础构件。在机械加工、热处理与装配过程中，齿面极易受到工具划擦或零件间外力磕碰，产生表面微细划痕与局部齿部缺损（主要表现为齿侧碰伤与齿角崩缺），这类缺陷容易在高速运转下诱发剧烈啮合振动与局部应力集中。传统车间主要依靠质检工人手持放大镜进行人工目检，工人劳动强度大，长时间连续作业极易产生视觉疲劳导致漏检，且不同检验员之间的评判尺度存在较强的主观性差异。

本项目源于我们在参观工厂生产车间时，目睹了质检工人肉眼排查齿面瑕疵的繁重工作，深切体会到采用自动化机器视觉检测替代人工肉眼质检的迫切性与工程价值。结合团队此前在 Python 编程、OpenCV 图像处理、PyTorch 深度学习以及 STM32 嵌入式开发上积累的实践基础，我们确立了“工业零件缺陷检测系统”这一实训研发方向。需要说明的是，系统现阶段明确界定为\textbf{辅助人工复核的筛查工具}，旨在流水线上自动初筛出疑似缺陷件供质检员重点复核，有效减轻人工肉眼负担。

\section*{二、 项目成员组成、特长、分工及协作配合与导师情况}

团队由 3 名大三工科学生组成跨方向协作小组：
\begin{itemize}
    \item \textbf{队长夏秋桐}：统筹规划总体进度，推进成员分工协作，负责 Jetson NX 边缘端部署、底层 STM32 微控制器编程、步进电机梯形平滑加减速控制算法编写
    \item \textbf{队员胡笑歌}：主要负责硬件设计版块，包括底层 STM32 微控制器编程、硬件串口通信协议联调，整体硬件架构设计与调试工作
    \item \textbf{队员陆怡澄}：主要负责软件与算法开发，包括自建数据集构建与标注、6 大轮算法演进、多模型异构融合与 Web 监控界面的编写
\end{itemize}

团队内部职责明确、定期召开技术研讨会，各环节配合默契。

\section*{三、 项目的创新点与特色}

项目的创新点与特色主要体现在数据集构建，模型优化与工程实现三个层面：

\subsection*{1. 数据集构建：“近重复组”划分与标签体系}
针对立体齿轮在多角度旋转拍摄过程中极易导致同一工件高度相似照片跨集合渗透的问题，团队设计了基于 64 位感知哈希（dHash，汉明距离 $\le 2$ 聚类合并）、拍摄时间序号及长宽比推断视角的“近重复组”聚类方法，严格以“组”为最小单元进行分层抽样划分数据集，减弱了数据泄漏导致的性能虚高；在标注规范上，历经多轮迭代确立了涵盖正常（\texttt{normal}）、微细划痕（\texttt{scratch}）、三种细分缺损位置（\texttt{bottom} 齿根、\texttt{oblique} 齿面、\texttt{side} 齿侧）以及困难样本（\texttt{difficult=1}，在训练集中予以剔除以防模糊边界导致梯度震荡）的 6 种严谨标签体系，保障了模型收敛稳定性与评测客观性。

\subsection*{2. 模型优化：分缺陷建模与多模型异构融合}
微细划痕（呈线状低对比度分布）与齿部缺损（呈块状局部几何形态破坏）在图像物理表征上形态差异极大。项目放弃了主流的通用单一模型检测方法，采用分缺陷分别建模策略：划痕子系统在 YOLO 目标检测主干中引入专门保留浅层高分辨率特征的 P2 检测层，并融合 EfficientNet-B0 与 ResNet18 图像分类网络；缺损子系统则结合单类 YOLO 与分类模型重点捕获形态突变。

针对异构模型输出置信度尺度不一的难题，项目在模型融合前通过验证集样本最小化负对数似然（NLL）进行温度缩放校准，独立学习各模型温度参数 $T$（如划痕子系统三成员分别校准为 0.700、2.050、2.125）；随后利用基于树状 Parzen 估计器（TPE）的贝叶斯优化算法在验证集上全局寻优确定最优集成权重（划痕子系统赋予目标检测器 0.75 高权重，缺损子系统均衡分配 0.50 权重）；最终在决策层采用“逻辑或规则（OR Rule）”合并两路报警，有效实现了分项专精与联合兜底。

\subsection*{3. 工程实现：软硬件闭环的全链路设计}
项目面向车间流水线在线初筛的实际工况，设计了“输送带逐件进料 $\rightarrow$ 透明亚克力载台无遮挡采光 $\rightarrow$ 上下双工业相机多视角覆盖 $\rightarrow$ Jetson Orin NX 本地推理 $\rightarrow$ STM32 控制步进电机平滑加减速分拣”的软硬件协同流水线。在科学价值上，深入探索了金属工业工件微弱缺陷在受控光照、多尺度特征提取与异构多模型融合下的识别规律；在工程应用价值上，作为一套低成本、轻量化的在线辅助筛查工具，能够显著降低质检工人的重复肉眼劳动强度。

\section*{四、 项目实施的进展情况及已取得的创新成果}

经过前期的扎实实验探索与软硬件工程开发，项目已圆满完成从实验室算法设计到第一代台架实物联调的阶段性任务：

\subsection*{1. 具体实施进度与技术演进历程}
\begin{itemize}
    \item \textbf{自主数据建库与 2 级流水线确立}：自行采购与定制 50 个真实齿轮，采集多视角图像并以“近重复组”防泄漏机制严格划分；确立了“ROI 快速定位外扩 4\% 裁剪 + 局部缺陷识别”的 2 级级联检测流水线，排除了整帧背景杂光干扰。
    \item \textbf{算法模型的 6 大轮演进探索}：
    \begin{enumerate}
        \item 第 1 轮尝试 Canny 边缘检测与自适应阈值分割等传统方法，受金属强反光亮斑及齿廓伪边缘干扰严重失效；
        \item 第 2 轮在整图直接训练通用 YOLO，微细划痕与车削刀纹高度混淆，漏检严重；
        \item 第 3 轮引入 ROI 裁剪二级流水线，对比了分类网络（整图粗判）与检测模型（框级定位）的指标差异；
        \item 第 4 轮尝试统一二分类大模型，微细划痕漏检依旧，盲目增加融合模型徒增计算开销；
        \item 第 5 轮确立分缺陷分别建模路线，分别构建划痕子系统与缺齿子系统；
        \item 第 6 轮确立温度缩放概率校准、TPE 贝叶斯权重寻优、经验奈曼–皮尔逊准则（FPR $\le 0.20$ 约束下选阈值）以及决策级“或规则”联合合并方案。
    \end{enumerate}
    \item \textbf{第一代整机实物台架搭建与闭环调试}：设计装配了平皮带输送机构与透明亚克力检测工位托盘载台；安装了上下双工业相机与白色漫射光源；开发了 STM32 步进电机梯形平滑加减速算法防止甩齿，定义了硬件串口（UART）指令；部署了本地 Jetson Orin NX 边缘端并完成定位模型 TensorRT FP16 编译加速，打通了“送料 $\rightarrow$ 采集 $\rightarrow$ 本地推理 $\rightarrow$ Web 监控 $\rightarrow$ 串口通信 $\rightarrow$ STM32 旋转分拣”的全物理闭环。
\end{itemize}

\subsection*{2. 已取得的阶段性创新成果}
\begin{itemize}
    \item \textbf{离线测试集检出能力显著改善}：在严格锁定的 150 张离线测试集（31 划痕、44 缺损、4 兼具，共 71 真实缺陷、79 正常）上，最终第 6 轮方案检出 63/71 张缺陷图，召回率达到 88.73\%；误报 18/79 张正常图，误报率为 22.78\%。相比早期第 4 轮统一分类基线（检出 54/71、误报 33/79），在相同数据下正确检出数净增 9 张，误报数减少 15 张，误报率从 41.77\% 下降至 22.78\%，验证了分缺陷建模、概率校准与贝叶斯权重寻优的综合效能。
    \item \textbf{实现了边缘端加速与成套可视化监控软件开发}：将齿轮定位模型成功编译为 TensorRT FP16 引擎，在 Orin NX 上的单环节耗时降至 26.7 ms，降低了系统端到端时延；上位机开发了轻量级 Web 监控界面，现场实时呈现多路图像、红框标记缺陷坐标、置信度得分及 PASS/REJECT 运行状态。
\end{itemize}
